\chapter[Women in History's Margins]{Why Women Were Written into History's Margins}
\section{A Genealogy with Only Half the Names}
First, pick up an object. Not an antiquity or a monument, but something tucked into an attic or the bottom of a chest in many Chinese homes: a family genealogy.

Open an old genealogy. If your family has none, try a library's local-history collection. You will find a peculiar layout: page after orderly page of names. An honored ancestor, his personal name, courtesy name, birth, death, burial place, and beneath him his sons. From the Ming dynasty to the Republican era, centuries branch out like a tree, without a single missing name.

Now try an experiment: find the women in this tree.

They are present, but differently. Beside a man's name appears "wife, of the X family" or "late mother, of the X family": this gentleman married a daughter of that household. Her surname is recorded; her personal name is not. The number of sons she bore is recorded, because they are new branches of the tree. Her birth year is often absent; her life receives not a word. Hundreds of pages of family history contain thousands of fully named men and hundreds of women with surnames but no personal names.

Try a closer experiment, without opening a book. Can you name your maternal great-grandmother---your mother's mother's mother? Count back just four generations, less than a century, and many people get stuck. Not because you do not care about your family, but because nobody thought this name needed to be solemnly handed down. Men's names entered the genealogy; women's hung in oral memory, a thread snapped within three generations.

This is how half of humanity appears in most historical records: present but unnamed; working but unrecorded; giving birth to everyone, yet not written into history.

This chapter asks why. Why were women long written into history's margins? Did they really do nothing, or did we keep accounts of only half of what happened?

\section{The Day the Census Taker Knocked}
Come with me to a scene reconstructed from the institutional details of its period. The people are fictional, but the forms, questions, and everything the pen recorded really happened.

Time: a Monday morning in early April 1891. Place: a Yorkshire town in northern England that grew around a coal mine. Coal smoke fills the air; sheets on washing lines never become entirely white.

Britain is conducting its census, a household-by-household count every ten years. Mr. Harper, the census taker, carries his forms to number seventeen, Miners' Lane. Tom, the man of the house, answers. He is thirty-four, newly home from the night shift, with dark circles around his eyes.

The form is completed by household. Harper asks the head's occupation. "Miner." He writes: miner.

And the wife? Her name, age, occupation?

Tom thinks a moment. "Sarah, thirty-two. She doesn't work. She's at home."

Under "occupation," Harper writes a word common on forms of the time: none.

Pause with that pen for a second. When did this day begin for Sarah, who "doesn't work"?

At half past four, she rises in darkness, lights the fire, and heats yesterday's water to fill a flask for her eldest son, who works the early shift. At five she fetches water from the well at the lane's end: four return trips, bruises on her arms never quite fading. At six she sets dough to rise, bakes bread, wakes six children, and changes the youngest. Later that morning she soaks neighbors' sheets in a large wooden tub, soapy water so cold her knuckles whiten. Taking in washing supplements many women's household income on this street, though no ledger calls it a business. A paying lodger also lives with the family; she cooks his three meals and makes his bed. Before noon she digs and sows the vegetable plot that will feed the family in summer. In the afternoon she darns socks, pickles vegetables, nurses a feverish little daughter, and watches the beans simmering for everyone. When Tom comes home that evening, hot water is ready.

Sarah has worked longer hours than her husband. Her labor sustains the household and brings in cash. Yet in the British Empire's official record, in the pile of forms guiding plans for schools, factories, and railways, she is a "none": a statistical zero.

Mr. Harper has not lied, and Tom means no harm. They simply follow a rule everyone takes for granted: work means leaving home, earning wages, and doing something that fits the form. Whatever happens inside the home, however exhausting or indispensable, does not count.

Remember the pen that wrote "none." The rest of this chapter asks why it writes that way.

\section{Busy All Day, Yet "Not Working"}
Let us lay out the conflict before answering it.

There are two apparently confident ways to read Sarah's day.

First: this is simply a division of labor. Tom spends nine hours underground digging coal, always at risk of collapse, flooding, or a gas explosion. Sarah manages the home, safely, while caring for the children. One works outside, the other inside: two hands, neither nobler than the other. The form merely counts "paid positions"; it does not declare housework unimportant.

Second: the problem is not division of labor but bookkeeping. Sarah's labor produces things, has value, and cannot be replaced. If she is ill for three days, the household immediately stalls, Tom must take leave, and the mine loses a worker. A system recording this labor as "none" is not a neutral ruler. It quietly defines "work," "contribution," and "importance." Those excluded by its definitions happen mainly to be women.

To see the dispute clearly, we need this chapter's central distinction: the \textbf{public sphere} and the \textbf{private sphere}. The public sphere comprises spaces "outside the home": markets, government, armies, courts, schools. Their activities have wages, contracts, archives, and monuments. The private sphere lies indoors: cooking, washing, childbirth, child-rearing, care for older and sick people. Here there are no wages, contracts, or archives. Work well done is treated as duty; it attracts attention only when nobody does it.

Society thus cuts life in half: one half counts as "serious affairs," the other as "getting through daily life." The crucial next step is that histories, statistical reports, monuments, and archives come almost entirely from the public sphere. Labor in the private sphere, together with its workers, consequently vanishes from the page. It was done, but not recorded; it mattered, but the definition of "important" had no box for it.

Now the real question emerges:

\textbf{Does history omit women's labor because it is objectively secondary, or do historical and statistical accounts first decide what counts as an "event," making women's labor appear never to have happened?}

In other words, which came first: the division of labor or the ledger?

This question concerns more than Sarah in 1891. It shapes today's debates over whether stay-at-home motherhood counts as work, why GDP cannot measure society's real wealth, and how many named women appear when you open a history textbook.

\section[The Case for a Woman's Place at Home]{"A Woman's Place Is in the Home": State the Case Fully}
Generations have argued over that question. Before going further, let us practice something this book repeatedly asks: give the position you disagree with its strongest reasons. This time it is the claim that "men naturally belong outside and women inside." The words may grate today, but unless you understand why they once persuaded so many people, you cannot understand their long reign.

Fully stated, the argument has at least four layers.

First, survival. Without infant formula, washing machines, or running water, feeding babies tethered mothers to home, while pregnancy and childbirth were themselves life-risking labor. Having one person stay nearby with children, gardens, and food preservation while another traveled to hunt, mine, or fight was not a conspiracy but a survival calculation. Second, protection. Mines collapse, battlefields kill, ships capsize. Why call excluding women from deadly occupations oppression? It leaves the greatest danger to men. Third, power exists inside families too. A Chinese mother-in-law's authority over daughters-in-law and household spending is real; an able housewife controls food, accounts, and social relations. Different spheres need not imply lower status: "different but equal." Fourth, stability. This division has operated for millennia, and the family is humanity's last refuge. However terrible the world outside, its light remains steady. Disturbing the division shakes civilization's foundations.

Each layer grasps something real: reproductive constraints, dangerous mines, domestic authority, the human need for stable order. That helps explain why men were not the system's only defenders. Ban Zhao of the Eastern Han was an accomplished woman who continued a historical work and taught imperial women at court, yet wrote Lessons for Women. Its opening chapter, "Humility and Weakness," broadly urged women to make modest submission their foundation. A system's defenders are never exclusively male; we will return to that.

The other side's voice is equally real, and cracks run through those four reasons.

Why is protection always one-way? Women "protected" from mines were also protected from wages, unions, and votes. Protection and prohibition often used the very same door. Why did domestic power so often depend on men's approval? A mother-in-law's authority often came after her husband's death, exercised provisionally as a "family elder." A widow's brief emergence into public view reveals that the position ordinarily did not belong to women. Most damagingly, if women's place in the family were natural, it should look broadly similar everywhere. In fact, societies arranged women's lives so differently that "nature" cannot contain the variation.

Consider several scenes.

More than three thousand years ago, ancient Egyptian women could own, inherit, buy, and sell land and houses, testify in court, and initiate divorce. Surviving contracts record these rights. More than thirty-eight hundred years ago, Babylon's Code of Hammurabi devoted dozens of provisions to the rights and duties of wives, widows, priestesses, and female tavern keepers. Women's place was among written law's earliest concerns, but its answers differed sharply from Egypt's.

At Japan's court around the year 1000, men wrote official documents and formal prose in classical Chinese, treating it as "proper writing." Kana was called "women's hand," a feminine script unfit for elevated use. What happened? The court lady Murasaki Shikibu used this disparaged "women's domain" to write The Tale of Genji, now regarded as one of the world's earliest novels. The small enclosure to which women were relegated eventually contained an era's highest literary achievement.

In nineteenth-century West Africa, the kingdom of Dahomey, in present-day Benin, maintained an all-female combat force among its military elite. Astonished European visitors called them the "Dahomey Amazons." Centuries earlier in North Africa, Fatima al-Fihri, a wealthy merchant's daughter in Fez, Morocco, used her inheritance to establish a mosque and promote learning. Al-Qarawiyyin, founded in 859, is traditionally regarded as one of the world's oldest continuously operating universities.

Women's circumstances also differed sharply within a single society. In classical Athens, wives of free citizens were confined to inner domestic rooms and rarely appeared publicly. Yet the city's many enslaved women performed every kind of "outside" work in markets, workshops, and other households. Elite women's seclusion rested on the labor of other women. "Women's position" has never been a single number: region, class, and ethnicity divide it into entirely different worlds.

The greatest weakness of "women naturally belong in the home" is therefore not its cruelty---in some periods it genuinely provided protection---but its failure to explain variation. Natural constants do not change from place to place; human-made variables do. Every society draws its own line around women's place. That means the line can be redrawn.

\section{A Bridge for Thought: Making Invisible Labor Visible}
Can we analyze a woman who "worked all day but has no work" with a portable tool rather than indignation alone? Yes: three simple questions. Use them to sift any historical narrative, set of statistics, or book about "great figures."

\textbf{First: who recorded this?}

Records do not fall from the sky. They have authors, and authors occupy positions: court historians are paid by the court, monks maintained by temples, reporters sent by newspapers, census forms designed by governments. The recorder's position determines where the light falls; everything else stays dark. Asking "who recorded this?" does not accuse anyone of lying. Most recorders are honest. It asks for whom, and for what purpose, the recording apparatus was built. An apparatus designed for taxation and conscription naturally sees able-bodied men clearly and misses the people spinning thread.

\textbf{Second: which box did it enter?}

Every record first divides the world into boxes, then fills them. A census has an "occupation" field, histories have criteria for admission to "biographies," national-income statistics define "production." The boxes speak more plainly than their contents: they directly reveal what an institution thinks worth counting. Remember the concept of \textbf{unpaid labor}: work receiving no wage and not passing through a market. Cooking at home is unpaid; restaurant cooking is paid. Looking after your child is unpaid; nursery care is paid. The same labor meets opposite fates in the accounts depending on whether money changes hands. Worldwide, women have long borne most unpaid labor. Asking "which box?" places the boundary between what counts and what does not on the table for inspection.

\textbf{Third: what counts as "important"?}

Even recorded events must be ranked. A battle, treaty, or dynastic change earns a chapter; delivering babies, raising children, saving seeds, and nursing the sick earn no words. Yet if nobody performed the latter for three days, the former "great events" could not happen. Nobody can fight a famous battle hungry and carrying plague. "Importance" is not an object's property but a license issued by a system of judgment.

Apply these three sieves to our opening objects. The genealogy: who recorded it? Male clan elders. Which box? "Descent," defined from father to son. What matters? Continuing the surname. Hence a woman who bore sons kept half a name, while one who did not evaporated altogether. The census: who recorded it? The state. Which box? "Occupation," defined as market employment. What matters? Taxable, plannable economic activity. Two cases, three questions: all find their mark.

This tool was not built only for women. Peasants, enslaved people, minority groups, children---everyone who seems to have "done little" in official histories can be examined with these questions. Often, "they left no records" really means "the recording apparatus was not built for them."

\section{"Remember the Ladies": A Letter Laughed Aside}
Now read a short primary document.

In March 1776, the North American colonies were moving toward independence. From their Massachusetts home, Abigail Adams wrote to her husband John, attending the Continental Congress in Philadelphia. As John and others drafted a new nation's blueprint, Abigail reminded him of something. Her phrase became one of American history's most famous lines from a letter:

\begin{quote}
"Remember the Ladies." ---Abigail Adams to John Adams, March 31, 1776
\end{quote}
Her message was broadly this: in the new code of laws, be more generous to women than your ancestors; do not place unlimited power in husbands' hands. Half seriously, she warned that if women received no place, they might organize a rebellion.

John's reply survives too. His response amounted to: your "extraordinary code" makes me laugh; do not take this seriously. He did laugh. Months later, the Declaration of Independence proclaimed everyone born equal, using the words "all men." Nationwide voting rights for American women came 144 years later. And, as we will see, many women waited several more decades for that written promise to become real.

Seventy-two years passed. In 1848, women and some men gathered in Seneca Falls, New York, for America's first women's-rights convention. Their Declaration of Sentiments made a clever move: it followed the Declaration of Independence's phrasing, changing one crucial point:

\begin{quote}
"We hold these truths to be self-evident: all men and women are born equal." ---Declaration of Sentiments, Seneca Falls, 1848
\end{quote}
Consider the strategy. They did not say the Declaration of Independence was wrong. They said: you are right, so live up to it. Use your opponent's text to expand your opponent's promise. This maneuver recurs throughout intellectual history; abolitionists, civil-rights activists, and countless successors also used it.

From Seneca Falls, movements for women's rights unfolded in successive surges, often called "waves" by historians.

\textbf{The first wave}, roughly the mid-nineteenth to early twentieth centuries, sought entry tickets to the public sphere: education, married women's property rights, and the vote. Victories came country by country. In 1893, New Zealand became the first country to grant women nationwide voting rights. Britain enfranchised property-qualified women over thirty in 1918, waiting until 1928 for equal terms with men. The United States passed the Nineteenth Amendment in 1920. French women did not vote for the first time until 1944. Switzerland, renowned for direct democracy, approved women's federal voting rights only in 1971, two years after humans landed on the Moon.

But the first wave was not united. America's fracture was particularly glaring: to secure Southern support, White suffrage organizations repeatedly relegated Black women to the back of marches or asked them not to attend. After the Nineteenth Amendment, Black women possessed voting rights on paper, yet Southern literacy tests, poll taxes, and violence kept them from polling stations for over forty more years, until the civil-rights movement of the 1960s. The indispensable lesson is that \textbf{"women" are not a naturally unified whole}. When race and class cut across them, "all women's shared condition" splits into very different conditions.

\textbf{The second wave} came roughly in the 1960s and 1970s. In The Second Sex (1949), French thinker Simone de Beauvoir famously argued, in essence, that one is not born a "woman" but gradually taught by society to become one. Following this insight, the movement pushed its battle lines inside the home. Why should housework belong to women by default? Could women decide about childbirth and contraception? Were marital violence and property "private matters" or political issues? Its slogan meant that the personal is political: supposedly individual family troubles often have institutions behind them.

The second wave also argued fiercely, including over a question still disputed today. One side sought "equality": exactly the same opportunities, tracks, and standards as men. Another asked why men's standards should be the standards. Care, nurture, and maintaining relationships deserve revaluation and serious social reward; women should not have to abandon them and squeeze onto a male-defined track to count as "liberated." Equality or revaluation? Stop at "the same as men," or rewrite what counts as valuable? There is no standard answer. Consider your own after reading this chapter.

China followed its own path, from Qiu Jin's poetry and execution to the New Culture Movement's attacks on foot-binding and arranged marriage. In 1950, the first law enacted after the founding of the People's Republic was the Marriage Law, placing the abolition of arranged marriage and equality between women and men at its forefront. The distance between paper and reality is another long story. It reminds us that changing a legal provision is easier than changing the ledger.

\section[Women's Marginalization: Three Accounts]{Why Were Women Marginalized? Three Explanations}
Return to the title's "why." Why did so many societies---not all societies---push women into a secondary position? Following this book's rule, distinguish what happened, a question of evidence, from why it happened, a question of explanation. Three competing accounts have their own territory and blind spots.

\textbf{First: bodies and technology.} Pregnancy and nursing are physiological burdens. In societies dependent on heavy plows, animal power, and warfare, institutions magnified differences in physical strength. Those able to guide the plow or wield the sword controlled key production and violence; power followed the distribution of muscles and wombs. This straightforward account explains the prevalence of gendered labor in agrarian societies. Its limits are clear too. Why was textile work, requiring less strength and largely performed by women, also undervalued? Why did female traders long dominate West African markets while free Athenian women were barred from appearing in them? More importantly, it smuggles a value judgment into "nature": why is a plow "higher" than a spindle? Bodies did not say so; people did.

\textbf{Second: property and inheritance.} This is Friedrich Engels's classic account in The Origin of the Family, Private Property and the State (1884). Broadly, once surplus and private property emerged, men wishing to leave possessions to children "certainly their own" had to control women's marriage and fidelity. Patrilineal families, chastity rules, and exclusion from public affairs consequently hardened. He called this "the world-historic defeat of the female sex." Its strength is explaining why very different societies repeatedly adopted similar arrangements: inheritance acted as a mold, stamping wherever it went. Its limitation is that later anthropology and archaeology found gender hierarchies even without private property, while evidence for a former golden age of women's rule is lacking. Real history is much messier than a single evolutionary line.

\textbf{Third: power over recording and definition.} This account digs deeper: marginalization concerns not only "what happened" but "what counted as having happened." Writing, legislating, ritual, historical compilation, newspaper publishing, and census design long rested mainly in men's hands. Male experience became "human" experience; male concerns became "great events." More than two thousand years ago, the Yang Huo chapter of the Analects stated:

\begin{quote}
"Only women and petty people are difficult to deal with." ---Analects, Yang Huo
\end{quote}
Commentators have argued for centuries over its targets and force, without consensus. The dispute itself shows that even the most authoritative text left only an ambiguous silhouette when mentioning women. This explanation illuminates why we know so little. Its limitation is that it explains marginalization's reproduction, not inequality's origins. The monopoly on writing itself still requires the first two explanations.

These accounts are not mutually exclusive alternatives but three links in a chain. Bodies and technologies supplied material for a division of labor; property and inheritance cast that division into hierarchy; control over writing and definitions made hierarchy invisible and self-reproducing, until it seemed "always to have been so."

Before leaving this section, consider a counterexample to a common misreading: "History progresses, so women's position naturally keeps improving." Classical Athens was Greek civilization's brilliant "golden age," a high point of philosophy, drama, and democracy. Yet free Athenian women were then confined to the home more strictly than in the earlier Homeric age. Civilizational "progress" and women's circumstances do not automatically move together. Sometimes an era's brilliance rests on locking half its people more tightly away. Whenever history becomes a straight line of "better every day," ask: better for whom?

\section[Gender Changes the View of History]{A Deeper Challenge: Gender Changes the View of History}
We now introduce this chapter's weightiest theoretical tool.

First distinguish two concepts. Physical characteristics at birth are called \textbf{biological sex}. The socially trained roles of "how boys and girls should behave"---clothing, studies, permission to speak, forbidden emotional displays---are called \textbf{gender}. The former comes largely from the body, the latter almost entirely from culture. Beauvoir's famous statement points to this distinction: womanhood is made, not simply born.

In 1986, historian Joan Scott published an influential article, "Gender: A Useful Category of Historical Analysis." Her proposal was far more radical than "write more women's history." Gender, she argued, is not merely for studying women; it is a lens for all history, because power loves gendered language when expressing and justifying itself. Conquerors describe conquered peoples as "weak, passive, feminine" to make domination seem appropriate. Employers label jobs "suitable for meticulous women" to make low pay seem reasonable. Gender is a language that colors power. Ask of any historical episode: how are "men" and "women" defined here, and whose interests does that definition serve?

Put on this lens, and three familiar kinds of history immediately look different.

\textbf{Military history.} Traditional military history covers generals, battles, and treaties: a panorama of the public sphere. A gender perspective first asks what armies eat, who grows it, who carries the wounded, and who sews uniforms. Women fill the answers, yet the "home front" remains an unnamed blank in older histories. Consider the First World War: casualties numbered in the tens of millions, men were drawn en masse into trenches, and millions of women suddenly filled factories, trams, hospitals, and fields. Nobody generously bestowed this opportunity; the war machine had no choice. Once fighting stopped, most countries promptly sent them back to the kitchen. But a door once opened could not close completely: women in Britain, Germany, and America obtained votes within a few years of the war. Some promises generated by mobilization were redeemed afterward. "Total war" inadvertently became a midwife to suffrage. Older military history would not mention this causal chain.

\textbf{Economic history.} We measure economies with ledgers such as GDP that recognize only market transactions. As discussed, unpaid labor counts as zero. The absurd result: paying a nanny to cook is "economic activity," while a mother cooking is "not economic." From the 1970s, economists and feminists began assigning market equivalents to housework and care. Methods and estimates vary widely, but even conservative calculations put them above a tenth of national economic output: enormous wealth long hidden by accounting. The Industrial Revolution also needs rewriting. British cotton mills, the workshops of Lowell, Massachusetts, and later Shanghai's textile factories relied chiefly on young rural women. Women supplied half industrialization's fuel, yet textbook illustrations often give the working class a male face by default. This field is moving from margin to center: the 2023 Nobel economics prize went to American economist Claudia Goldin for research on women's labor markets.

\textbf{Political history.} Traditional political history counts monarchs, presidents, and laws. A gender lens reminds us that the electorate's boundaries are themselves products of political struggle. In most countries, the answer to "who may vote?" still read "men" well into the twentieth century. Expanding suffrage also changes political content. Laws on maternity leave, childcare, marital property, and domestic violence reached agendas close to the time women gained votes and entered parliaments. Institutionally, "the personal is political" meant pulling life "inside the home" out of legislation's blind spot.

Weigh both this theory's power and its limits. A gender perspective illuminates huge blind spots, but it is not the only light. Not every hierarchy reduces to gender. Translating everything into "men and women" misses equally durable structures of class and ethnicity. Remember Black women being asked to stand at the back of the Seneca Falls procession. A good tool does not answer everything for you; it makes the old pretense of not seeing impossible.

\section{The Line Is Still Here Today}
This century-old question is beside you now, although the forms have changed.

Today's census taker will no longer write "unemployed" beside your mother's name. But the old form has descendants. Time-use surveys published by China's National Bureau of Statistics show women spending substantially more time on unpaid housework and care than men, roughly more than twice as much. Sarah's wooden tub still holds water; it has simply become the position beside a dishwasher. The question has shifted from "does it count as work?" to "why is it hers by default?"

At school, the line speaks a different dialect. "Girls are good at humanities; boys have more long-term potential." You may have heard it. It assumes boys' present setbacks are temporary and girls' present lead is accidental: identical marks, different stories. The line restrains boys too. Male nursing students and preschool teachers face another inspecting gaze. This is gender's actual shape: not a cage designed exclusively for one sex, but a script assigned to each, with punishment for refusing the role.

At home and online, its loose ends appear everywhere. A child's surname can ignite a trending dispute. "Woman driver" is treated as shorthand for incompetence, though traffic-accident statistics do not support the joke. Employers' questions about female applicants' marriage and childbirth plans are already violations repeatedly prohibited by labor authorities. Our opening experiment still works: count the named women across ten units of your history textbook. You probably will not need all ten fingers.

Following this book's rule, distinguish four kinds of statement. Verified facts: unpaid labor is unequally distributed, and textbooks and historical records contain large gender imbalances. A reasonable explanation: the chain of labor division, institutions, and bookkeeping continues moving under its own inertia; nobody was simply born this way. The author's judgment: inertia means continued motion unless force intervenes. The open question: where should that force be applied? That is our final section's question for you.

\section[Who Decides What Counts as Work?]{Your Question: Who Decides What Counts as Work?}
Return to the genealogy and the census taker's pen.

Now answer a question without a standard answer, and give reasons:

\textbf{If, starting tomorrow, cooking, laundry, childcare, and eldercare all received explicit prices, wages, and entries in national economic output, would the world be fairer?}

Before answering, weigh both sides.

You already have reasons for pricing. Recognition is fairness's first step. Sarah's labor was recorded as "none" not because she did little but because the ledger could not see it. Change the accounts, and stay-at-home parents need not leave divorce empty-handed; elder caregivers gain security in old age; "you don't earn anything" loses its sting. What remains invisible cannot be protected. Priced labor can enter the law.

The case against pricing is also substantial. A price recognizes, but may also absorb. What is a mother's hug worth? Once everything passes over the market's scales, nowhere remains for doing things "not for money." Home is home precisely because accounts are not kept there. Handing the entire private sphere over to public-sphere logic may not remove the boundary; it may let the public sphere swallow the private.

Push further: is there a third way, neither pretending housework does not exist nor turning it into a business? Can society "record" such labor differently?

Finally, return to the genealogy. Tonight, ask your mother's mother what her mother was called. If nobody can answer, consider whether this "don't know" and the "none" on the 1891 census were written by the same pen. With which generation should we begin writing differently?

There is no standard answer. But from today, you have three questions, a conceptual distinction, and a lens capable of rewriting your whole view of history. The tools are yours. It is your turn to draw the line.
